% ===== tr-03 块5：第 3 章 pp.36--38（原文页 36--38）=====
% [原文 p.36]
这里有一个技术问题：单性的定义假定存在环面纤维化 $f:X\to B$。而我们希望完全用 $B$ 本身来定义单性。不幸的是，这个解决方案相当技术化，所产生的定义非常难以消化（[22] 的定义 1.60）。这里我们只说：若 $B$ 在这个新意义下是简单的，且 $X(B_0)\to B_0$ 以某种合理的方式紧化为拓扑环面纤维化 $f:X(B)\to B$，则只要 $\dim B\le3$，$f$ 在第 1 节的意义下就是简单的。在高维情形，这个新的单性不一定蕴涵旧的单性；见 Helge Ruddat 即将完成的博士论文。这出现在 $X(B)$ 中产生轨形奇点的情形；众所周知，这类奇点在高维无法避免。

一旦接受这个定义，生活就大大简化。非同寻常的是：配定了对偶相交复形的对数 Calabi--Yau 空间的模空间，先验地非常复杂，但当 $B$ 简单时却有非常简单的描述！[22] 的一个非常困难的主要结果（定理 5.4）是

\begin{theorem}{3.8.8}
给定正的、简单的 $(B,\mathcal P)$，以 $(B,\mathcal P)$ 为对偶相交复形的对数 Calabi--Yau 空间之集，模去保持 $B$ 的同构，是 $H^1(B,i_*\Lambda\otimes k^\times)$。若一个同构在对偶相交复形上诱导恒等映射，则称它保持 $B$。
\end{theorem}

于是模空间是一个代数环面（或若干代数环面的不交并），其维数等于 $\dim_kH^1(B,i_*\Lambda\otimes k)$。

注意这正是 SYZ 猜想所预言的期望维数。事实上，若 $X_0(B,\mathcal P,s)^\dagger$ 的光滑化存在，且它是 $X_0(B)$ 的拓扑紧化 $X(B)$、带有延拓 $f_0:X(B_0)\to B_0$ 的简单环面纤维化 $f:X(B)\to B$，则 $R^{n-1}f_{0*}\R\cong\Lambda_\R$，故由单性，$R^{n-1}f_*\R\cong i_*\Lambda_\R$。第 1 节的讨论表明 $\dim H^1(B,R^{n-1}f_*\R)$ 是该光滑化的 $h^{1,n-1}$，而这当然就是该光滑化的复模空间的维数。
% [Tutor Note: 上式里的 $i_*\Lambda_\R$ 是局部系统 $R^{n-1}f_*\R$ 的记号：$i:B_0\hookrightarrow B$ 是开浸入，$\Lambda$ 是整仿射函数层；不要把它读成 $i$ 的上同调。另注意 $h^{1,n-1}$ 是 Hodge 数，不是 $h^{1}\cdot h^{n-1}$]

这一论证可以通过引入\emph{对数微分}而变得严格。

\begin{definition}{3.8.9}
设 $\pi:X^\dagger\to S^\dagger$ 是对数空间的态射。$X^\dagger$ 在 $S^\dagger$ 上、取值于 $\mathcal O_X$-模 $\mathcal E$ 的一个\emph{对数导子}，是一对 $(D,\mathrm{Dlog})$，其中 $D:\mathcal O_X\to\mathcal E$ 是 $X/S$ 的通常导子，$\mathrm{Dlog}:\mathcal M_X^{\mathrm{gp}}\to\mathcal E$ 是阿贝尔层同态且 $\mathrm{Dlog}\circ\pi^\#=0$；它们满足如下相容性条件
\[
D(\alpha_X(m))=\alpha_X(m)\cdot\mathrm{Dlog}(m),
\]
对一切 $m\in\mathcal M_X$ 成立。
\end{definition}

我们以 $\Theta_{X^\dagger/S^\dagger}$ 记 $X^\dagger$ 在 $S^\dagger$ 上、取值于 $\mathcal O_X$ 的对数导子层。置 $\Omega^1_{X^\dagger/S^\dagger}=\operatorname{Hom}_{\mathcal O_X}(\Theta_{X^\dagger/S^\dagger},\mathcal O_X)$。

这推广了更熟悉的概念：沿正规交叉除子的、具有对数极点的微分。若 $Y\subseteq X$ 是正规交叉除子、在 $X$ 上诱导一个对数结构，则 $\Omega^1_{X^\dagger/k}$ 是沿 $Y$ 具有对数极点的微分层，而 $\Omega^1_{Y^\dagger/k}$ 是此层到 $Y$ 的限制。一般地，若 $\pi$ 对数光滑，则 $\Omega^1_{X^\dagger/S^\dagger}$ 局部自由。因此，可以在对数光滑态射的对数范畴中做形变理论
% [原文 p.37]
（见 [34]）。这是我们在这幅图景中引入对数几何的主要原因之一。

若 $X_0(B,\mathcal P,s)^\dagger\to0^\dagger$ 是对数 Calabi--Yau 空间，则该态射在 $Z$ 之外对数光滑。定义
\[
\begin{aligned}
\Theta^p_{X_0(B,\mathcal P,s)}&:=j_*\bigwedge^p\Theta_{(X_0(B,\mathcal P,s)^\dagger\setminus Z)/0^\dagger},\\
\Omega^p_{X_0(B,\mathcal P,s)}&:=j_*\bigwedge^p\Omega^1_{(X_0(B,\mathcal P,s)^\dagger\setminus Z)/0^\dagger},
\end{aligned}
\]
其中 $j:X_0(B,\mathcal P,s)\setminus Z\to X_0(B,\mathcal P,s)$ 是包含。

于是有

\begin{theorem}{3.8.10}
设 $(B,\mathcal P)$ 是正的、简单的，并设给定对数 Calabi--Yau 空间 $X_0(B,\mathcal P,s)^\dagger\to0^\dagger$，它作为某个环面退化 $\mathcal X\to D$ 的中心纤维出现，而该退化的一般纤维 $\mathcal X_t$ 非奇异。则对 $q=0,1,n-1$ 与 $n$（其中 $n=\dim B$），有同构
\[
\begin{aligned}
H^p(B,i_*\bigwedge\nolimits^q\Lambda\otimes k)&\cong H^p(X_0(B,\mathcal P,s),\Theta^q_{X_0(B,\mathcal P,s)})\cong H^p(\mathcal X_t,\Theta^q_{\mathcal X_t})\\
H^p(B,i_*\bigwedge\nolimits^q\check\Lambda\otimes k)&\cong H^p(X_0(B,\mathcal P,s),\Omega^q_{X_0(B,\mathcal P,s)})\cong H^p(\mathcal X_t,\Omega^q_{\mathcal X_t})
\end{aligned}
\]
其中 $\Theta^q_{\mathcal X_t}$ 与 $\Omega^q_{\mathcal X_t}$ 是光滑纤维 $\mathcal X_t$ 上全纯多向量场与全纯微分的通常层。
\end{theorem}

这一结果的证明连同若干其他结果出现在 [24] 中。在附加假设下（本质上说 $\mathcal X_t$ 的镜像非奇异），该结果对一切 $q$ 成立。特别地注意：由于 $\Lambda$ 与 $\check\Lambda$ 在离散 Legendre 变换下互换，我们从这一结果得到通常 Hodge 数的互换。在出现的 Calabi--Yau 为奇异的更一般情形下，人们可以猜测这些群、实际 Hodge 数与弦 Hodge 数（stringy Hodge numbers）之间的关系。这些问题在 Helge Ruddat 即将完成的博士论文中讨论。
% [Tutor Note: 定理 3.8.10 的两行分别是 $\bigwedge^q\Lambda$（整仿射侧）与 $\bigwedge^q\check\Lambda$（对偶侧）的外幂；第二行右端用 $\Omega^q$（全纯微分）而非 $\Theta^q$（多向量场），两行不要抄串。另：smooth fibre 用英式拼法 fibre，故译"纤维"；stringy Hodge numbers 有译"弦 Hodge 数"者，此处保留英文原词于括号内]

\section{锥图景与扇图景}

本节纯属哲学性的讨论。我们在第 7、8 节的大部分讨论中聚焦于对偶相交复形，特别是聚焦于"从其对偶相交复形构造退化"这一问题。由于我们的首要目标是解决重构问题------问题 3.5.6，而对偶相交复形与复结构相关（定理 3.7.7），聚焦于对偶相交复形似乎很自然。我们将在下一节看到，这种直觉未必总是对的。到目前为止，相交复形似乎只在谈论镜像对称时才出现。然而，镜像对称教导我们把图景的两侧放在同等地位上。当我们由对偶相交复形构造退化 Calabi--Yau 空间时，我们说处于\emph{扇图景}；而当我们由相交复形构造退化 Calabi--Yau 空间时，我们说处于\emph{锥图景}。

更确切地说，我们已经看到：给定带环面多面体分解 $(B,\mathcal P)$ 的带奇点整仿射流形，再选定开粘合数据 $s$，就确定了一个空间 $X_0(B,\mathcal P,s)$，连同幺半群层 $\overline{\mathcal M}_{X_0(B,\mathcal P,s)}$。某些附加数据可以指定 $X_0(B,\mathcal P,s)$ 上以该幽灵层为幽灵层的对数结构。$X_0(B,\mathcal P,s)$ 的不可约分支是用扇来定义的------由 $\mathcal P$ 在 $\mathcal P$ 的每个顶点处的扇结构给出。这就是我们把这一侧称为扇图景的原因。
% [原文 p.38]
另一方面，给定 $(B,\mathcal P)$，在适当的粘合数据 $\check s$ 下我们还可以构造一个\emph{射影}概形 $\check X_0(B,\mathcal P,\check s)$。该概形的不可约分支与 $\mathcal P$ 的极大胞腔一一对应；给定这样的极大胞腔 $\sigma$，把它看作 $\R^n$ 中的格多面体便确定一个射影环面簇，而 $\check X_0(B,\mathcal P,\check s)$ 就是用数据 $\check s$ 把这些射影环面簇粘合起来而得到的。这与扇图景中出现的数据并不完全相同，因为我们还需要粘合线丛，而这是附加数据。把这一侧称为锥图景的理由是：每个不可约分支可以如下描述。给定 $\sigma\subseteq M_\R$，置 $P_\sigma=C(\sigma)\cap(M\oplus\Z)$。则相应的射影环面簇是 $\operatorname{Proj}\C[P_\sigma]$，其中 $\C[P_\sigma]$ 用 $P_\sigma$ 到 $\Z$ 的投影来分次。因此不可约分支与分层来自 $\mathcal P$ 中元素上的锥。我们把锥图景与扇图景之间的对偶小结如下：
% [Tutor Note: 末句原文写 cones over elements of 后面跟花体 P（多面体分解）；但按上下文与 $P_\sigma=C(\sigma)\cap(M\oplus\Z)$ 的用法，此处大概本意是幺半群 $P$。1200 dpi 裁图确认原文字形为花体，正文照原文排]
% [Tutor Note: 下面的表在原文中无表号、无表题，用全框线排版；本译本不自造表号。表中第 3 行"$\varphi$ 多值凸分段线性函数"对应扇侧的"丰沛线丛"与锥侧的"幺半群层"，方向不要弄反：扇侧的数据是 $\varphi$，锥侧的数据是 $\overline{\mathcal M}$]
\begin{center}
\renewcommand{\arraystretch}{1.3}
\resizebox{0.97\textwidth}{!}{%
\begin{tabular}{|l|l|l|}
\hline
 & 扇图景 & 锥图景\\
\hline
粘合数据给出 & $X_0(B,\mathcal P,s)$，$\overline{\mathcal M}_{X_0(B,\mathcal P,s)}$ & $\check X_0(B,\mathcal P,\check s)$，丰沛线丛\\
\hline
$\sigma\in\mathcal P,\ \dim\sigma=p$ & $X_0(B,\mathcal P,s)$ 的 $n-p$ 维分层 & $\check X_0(B,\mathcal P,\check s)$ 的 $p$ 维分层\\
\hline
$\varphi$：多值凸分段线性函数 & $X_0(B,\mathcal P,s)$ 上的丰沛线丛 & 幺半群层 $\overline{\mathcal M}_{\check X_0(B,\mathcal P,\check s)}$\\
\hline
$H^p(B,i_*\bigwedge\nolimits^q\Lambda\otimes k)$ & $H^p(X_0(B,\mathcal P,s),\Theta^q_{X_0(B,\mathcal P,s)})$ & $H^p(\check X_0(B,\mathcal P,\check s),\Omega^q_{\check X_0(B,\mathcal P,\check s)})$\\
\hline
\end{tabular}%
}
\end{center}

（两侧的粘合数据可能都需要一些限制，$\varphi$ 才能给出所需的数据。）

特别地，镜像对称把关于对数结构的离散信息（即 $\overline{\mathcal M}_{X_0(B,\mathcal P,s)}$）与关于极化的离散信息（即每个不可约分支上线丛的类）互换。

\section{热带曲线}

到目前为止，从这个代数几何版本的 SYZ 中我们只看到镜像对称最初步的方面，例如 Hodge 数的互换。然而，这一途径的真正趣味在于种种迹象：它将对镜像对称中的有理曲线计数给出自然的解释。若遵从上一节的哲学，我们需要在带奇点的仿射流形上辨认出这样的结构：它在两种图景之一中与有理曲线有关，而在另一种图景中与周期有关。我相信应当研究的是带奇点仿射流形 $B$ 上的热带曲线结构。关于 $\R^n$ 中热带曲线的介绍，见 [46]、[50]。这里，$B$ 可以取为热带仿射流形而非整仿射流形；"热带"之名即由此而来。
